% Part: normal-modal-logic
% Chapter: tableaux
% Section: completeness

\documentclass[../../../include/open-logic-section]{subfiles}

\begin{document}

\olfileid[pt]{nml}{tab}{cpl}

\olsection{Completude para \Log{K}}

\begin{explain}
  Para mostrar que o método dos !!{tableau}s é completo, precisamos
  mostrar que, sempre que não há um !!{tableau} fechado que mostre
  $\Gamma \Proves !A$, então $\Gamma \Entails/ !A$, O.s.\@, há um
  contramodelo. Mas ``não há um !!{tableau} fechado'' significa que toda
  maneira pela qual poderíamos tentar construir um deve deixar de fechar.
  O truque consiste em perceber que, se toda maneira desse tipo deixa de
  fechar, então uma maneira específica, \emph{sistemática e exaustiva},
  também deixa de fechar. E essa maneira sistemática e exaustiva fecharia
  se existisse um !!{tableau} fechado. O único tableau conterá, entre seus
  ramos abertos, toda a informação necessária para definir um contramodelo.
  O contramodelo fornecido por um ramo aberto desse tableau conterá como
  mundos todos os prefixos usados nesse ramo, e !!a{propositional
    variable}~$p$ será verdadeira em $\sigma$ sse
  $\sFmla{\True}{p}[\sigma]$ ocorrer no ramo.
\end{explain}

\begin{defn}
  Um ramo de !!a{tableau} é chamado completo se, sempre que contém uma
  !!{formula} prefixada $\sFmla{S}{!A}[\sigma]$ à qual uma regra pode
  ser aplicada, também contém:
  \begin{enumerate}
    \item as !!{formula}s prefixadas que são as conclusões correspondentes
      da regra, no caso das regras proposicionais de empilhamento;
    \item uma das !!{formula}s de conclusão correspondentes, no caso das
      regras proposicionais de ramificação;
    \item ao menos uma conclusão possível, no caso das regras modais que
      exigem um prefixo novo;
    \item a conclusão correspondente para todo prefixo que ocorre no ramo,
      no caso das regras modais que exigem um prefixo usado.
  \end{enumerate}
\end{defn}

\begin{explain}
Por exemplo, um ramo completo contém $\sFmla{\True}{!B}[\sigma]$ e
$\sFmla{\True}{!C}[\sigma]$ sempre que contém $\sFmla{\True}{!B
  \land !C}$. Se contém $\sFmla{\True}{!B \lor !C}[\sigma]$, contém
ao menos uma entre $\sFmla{\False}{!B}[\sigma]$ e
$\sFmla{\True}{!C}[\sigma]$. Se contém \iftag{prvBox}
{$\sFmla{\False}{\Box}[\sigma]$} {$\sFmla{\True}{\Diamond}[\sigma]$}
também contém \iftag{prvBox} {$\sFmla{\False}{\Box}[\sigma.n]$}
{$\sFmla{\True}{\Diamond}[\sigma.n]$} para ao menos um~$n$. E, sempre
que contém \iftag{prvBox} {$\sFmla{\True}{\Box}[\sigma]$}
{$\sFmla{\False}{\Diamond}[\sigma]$}, também contém \iftag{prvBox}
{$\sFmla{\True}{\Box}[\sigma.n]$}
{$\sFmla{\False}{\Diamond}[\sigma.n]$} para todo~$n$ tal que
$\sigma.n$ seja usado no ramo.
\end{explain}

\begin{prop}\ollabel{prop:complete-tableau}
  Todo $\Gamma$ finito tem !!a{tableau} no qual todo ramo é completo.
\end{prop}

\begin{proof}
  Considere um ramo aberto de !!a{tableau} para~$\Gamma$. Há um número
  finito de !!{formula}s prefixadas no ramo às quais uma regra poderia
  ser aplicada. Em alguma ordem fixa (digamos, de cima para baixo), para
  cada uma dessas !!{formula}s prefixadas para a qual as condições
  (1)--(4) ainda não valem, aplique as regras aplicáveis a ela para
  estender o ramo. Em alguns casos, isso produzirá ramificações; aplique
  a regra na extremidade de cada ramo resultante para todas as !!{formula}s
  prefixadas restantes. Como o número de !!{formula}s prefixadas é finito
  e o número de prefixos usados no ramo é finito, esse procedimento acaba
  produzindo (possivelmente muitos) ramos que estendem o ramo original.
  Aplique o procedimento a cada um e repita. Mas, por construção, todo
  ramo é fechado.
\end{proof}

\begin{thm}[Completude]
  \ollabel{thm:tableau-completeness}
  Se $\Gamma$ não tem !!{tableau} fechado, $\Gamma$ é satisfatível.
\end{thm}

\begin{proof}
Pela proposição, $\Gamma$ tem !!a{tableau} no qual todo ramo é completo.
Como não tem !!{tableau} fechado, tem !!a{tableau} no qual ao menos um
ramo é aberto e completo. Seja $\Delta$ o conjunto das !!{formula}s
prefixadas no ramo e seja $P(\Delta)$ o conjunto dos prefixos que nele
ocorrem.

Definimos um modelo~$\mModel{M(\Delta)} = \tuple{P(\Delta), R, V}$ no
qual os mundos são os prefixos que ocorrem em~$\Delta$, e a relação de
acessibilidade é dada por:
\[
R\sigma\sigma' \quad \text{sse} \quad
\sigma'=\sigma.n \quad \text{para algum~$n$}
\]
e
\[
V(p) = \Setabs{\sigma}{\sFmla{\True}{p}[\sigma] \in \Delta}.
\]
Mostramos por indução sobre~$!A$ que, se
$\sFmla{\True}{!A}[\sigma] \in \Delta$, então
$\mSat{M(\Delta)}{!A}[\sigma]$ e, se
$\sFmla{\False}{!A}[\sigma] \in \Delta$, então
$\mSat/{M(\Delta)}{!A}[\sigma]$.

\begin{enumerate}
  \item \indcase{!A}{p}{Se $\sFmla{\True}{\indfrm}[\sigma] \in \Delta$,
    então $\sigma \in V(p)$ (pela definição de~$V$) e, portanto,
    $\mSat{M(\Delta)}{\indfrm}[\sigma]$.

    Se $\sFmla{\False}{\indfrm}[\sigma] \in \Delta$, então
    $\sFmla{\True}{\indfrm}[\sigma] \notin \Delta$, pois, do contrário,
    o ramo seria fechado. Logo, $\sigma \notin V(p)$ e
    $\mSat/{M(\Delta)}{\indfrm}[\sigma]$.}
  \item \indcase{!A}{\lnot !B}{\iftag{probNot}{Exercício.}{Se
      $\sFmla{\True}{\indfrm}[\sigma] \in \Delta$, então
      $\sFmla{\False}{!B}[\sigma] \in \Delta$, pois o ramo é completo.
      Pela hipótese de indução, $\mSat/{M(\Delta)}{!B}[\sigma]$ e, portanto,
      $\mSat{M(\Delta)}{\indfrm}[\sigma]$.

      Se $\sFmla{\False}{\indfrm}[\sigma] \in \Delta$, então
      $\sFmla{\True}{!B}[\sigma] \in \Delta$, pois o ramo é completo.
      Pela hipótese de indução, $\mSat{M(\Delta)}{!B}[\sigma]$ e, portanto,
      $\mSat/{M(\Delta)}{\indfrm}[\sigma]$.}}
  \item \indcase{!A}{!B \land !C}{\iftag{probAnd}{Exercício.}{Se
      $\sFmla{\True}{\indfrm}[\sigma] \in \Delta$, então tanto
      $\sFmla{\True}{!B}[\sigma] \in \Delta$ quanto
      $\sFmla{\True}{!C}[\sigma] \in \Delta$, pois o ramo é completo.
      Pela hipótese de indução, $\mSat{M(\Delta)}{!B}[\sigma]$ e
      $\mSat{M(\Delta)}{!C}[\sigma]$. Portanto,
      $\mSat{M(\Delta)}{\indfrm}[\sigma]$.

      Se $\sFmla{\False}{\indfrm}[\sigma] \in \Delta$, então ou
      $\sFmla{\False}{!B}[\sigma] \in \Delta$, ou
      $\sFmla{\False}{!C}[\sigma] \in \Delta$, pois o ramo é completo.
      Pela hipótese de indução, ou $\mSat/{M(\Delta)}{!B}[\sigma]$, ou
      $\mSat/{M(\Delta)}{!B}[\sigma]$. Portanto,
      $\mSat/{M(\Delta)}{\indfrm}[\sigma]$.}}
  \item \indcase{!A}{!B \lor !C}{\iftag{probOr}{Exercício.}{Se
      $\sFmla{\True}{\indfrm}[\sigma] \in \Delta$, então ou
      $\sFmla{\True}{!B}[\sigma] \in \Delta$, ou
      $\sFmla{\True}{!C}[\sigma] \in \Delta$, pois o ramo é completo.
      Pela hipótese de indução, ou $\mSat{M(\Delta)}{!B}[\sigma]$, ou
      $\mSat{M(\Delta)}{!C}[\sigma]$. Portanto,
      $\mSat{M(\Delta)}{\indfrm}[\sigma]$.

      Se $\sFmla{\False}{\indfrm}[\sigma] \in \Delta$, então tanto
      $\sFmla{\False}{!B}[\sigma] \in \Delta$ quanto
      $\sFmla{\False}{!C}[\sigma] \in \Delta$, pois o ramo é completo.
      Pela hipótese de indução, tanto $\mSat/{M(\Delta)}{!B}[\sigma]$
      quanto $\mSat/{M(\Delta)}{!B}[\sigma]$. Portanto,
      $\mSat/{M(\Delta)}{\indfrm}[\sigma]$.}}
  \item \indcase{!A}{!B \lif !C}{\iftag{probIf}{Exercício.}{Se
      $\sFmla{\True}{\indfrm}[\sigma] \in \Delta$, então ou
      $\sFmla{\False}{!B}[\sigma] \in \Delta$, ou
      $\sFmla{\True}{!C}[\sigma] \in \Delta$, pois o ramo é completo.
      Pela hipótese de indução, ou $\mSat/{M(\Delta)}{!B}[\sigma]$, ou
      $\mSat{M(\Delta)}{!C}[\sigma]$. Portanto,
      $\mSat{M(\Delta)}{\indfrm}[\sigma]$.

      Se $\sFmla{\False}{\indfrm}[\sigma] \in \Delta$, então tanto
      $\sFmla{\True}{!B}[\sigma] \in \Delta$ quanto
      $\sFmla{\False}{!C}[\sigma] \in \Delta$, pois o ramo é completo.
      Pela hipótese de indução, tanto $\mSat{M(\Delta)}{!B}[\sigma]$
      quanto $\mSat/{M(\Delta)}{!B}[\sigma]$. Portanto,
      $\mSat/{M(\Delta)}{\indfrm}[\sigma]$.}}
  \item \indcase{!A}{\Box !B}{\iftag{probBox}{Exercício.}{Se
      $\sFmla{\True}{\indfrm}[\sigma] \in \Delta$, então, como o ramo é
      completo, $\sFmla{\True}{!B}[\sigma.n] \in \Delta$ para todo
      $\sigma.n$ usado no ramo, O.s.\@, para todo $\sigma' \in P(\Delta)$
      tal que $R\sigma\sigma'$. Pela hipótese de indução,
      $\mSat{M(\Delta)}{!B}[\sigma']$ para todo $\sigma'$ tal que
      $R\sigma\sigma'$. Portanto,
      $\mSat{M(\Delta)}{\indfrm}[\sigma]$.

      Se $\sFmla{\False}{\indfrm}[\sigma] \in \Delta$, então, para algum
      $\sigma.n$, $\sFmla{\False}{!B}[\sigma.n] \in \Delta$, pois o ramo
      é completo. Pela hipótese de indução,
      $\mSat/{M(\Delta)}{!B}[\sigma.n]$. Como $R\sigma(\sigma.n)$, há um
      ~$\sigma'$ tal que $\mSat/{M(\Delta)}{!B}[\sigma']$. Portanto,
      $\mSat/{M(\Delta)}{\indfrm}[\sigma]$.}}
  \item \indcase{!A}{\Diamond !B}{\iftag{probDiamond}{Exercício.}{Se
      $\sFmla{\True}{\indfrm}[\sigma] \in \Delta$, então, para algum
      $\sigma.n$, $\sFmla{\True}{!B}[\sigma.n] \in \Delta$, pois o ramo
      é completo. Pela hipótese de indução,
      $\mSat{M(\Delta)}{!B}[\sigma.n]$. Como $R\sigma(\sigma.n)$, há um
      ~$\sigma'$ tal que $\mSat{M(\Delta)}{!B}[\sigma']$. Portanto,
      $\mSat{M(\Delta)}{\indfrm}[\sigma]$.

      Se $\sFmla{\False}{\indfrm}[\sigma] \in \Delta$, então, como o
      ramo é completo, $\sFmla{\False}{!B}[\sigma.n] \in \Delta$ para
      todo $\sigma.n$ usado no ramo, O.s.\@, para todo
      $\sigma' \in P(\Delta)$ tal que $R\sigma\sigma'$. Pela hipótese de
      indução, $\mSat/{M(\Delta)}{!B}[\sigma']$ para todo $\sigma'$ tal
      que $R\sigma\sigma'$. Portanto,
      $\mSat/{M(\Delta)}{\indfrm}[\sigma]$.}}
\end{enumerate}
Como $\Gamma \subseteq \Delta$, $\mSat{M(\Delta)}{\Gamma}$.
\end{proof}

\begin{prob}
Complete a demonstração de \olref[nml][tab][cpl]{thm:tableau-completeness}.
\end{prob}

\begin{cor}
\ollabel{cor:entailment-completeness}
Se $\Gamma \Entails !A$, então $\Gamma \Proves !A$.
\end{cor}

\begin{cor}
\ollabel{cor:weak-completeness}
Se $!A$ é verdadeira em todos os modelos, então $\Proves !A$.
\end{cor}

\end{document}
